\documentclass[a4paper,11pt]{article}

\usepackage[a4paper, margin=2.5cm]{geometry}
\usepackage{hyperref}
\usepackage{parskip}
\usepackage{listings}
\usepackage{xcolor}
\usepackage{titlesec}
\usepackage{enumitem}

\hypersetup{
    colorlinks=true,
    linkcolor=blue,
    urlcolor=blue
}

\lstset{
    basicstyle=\ttfamily\small,
    backgroundcolor=\color{gray!10},
    frame=single,
    breaklines=true
}

\title{\textbf{Constructing Hopf--Galois structures and skew bracoids of small degree}\\[0.5em]
\large Enumeration algorithm --- guide to use}
\author{Andrew Darlington \and E.A.\ O'Brien}
\date{Last revised March 2026}

\begin{document}

\maketitle

This directory contains material related to the enumeration algorithm described in the preprint:

\begin{center}
\textit{``Constructing Hopf--Galois structures and skew bracoids of small degree''}\\
Authors: Andrew Darlington and E.A.\ O'Brien
\end{center}

\section*{1. Directory structure}

\begin{itemize}[noitemsep]
    \item \texttt{manager.m}, \texttt{worker.m}, and various shell scripts in the top-level directory.
    \item \texttt{computed-data/} --- contains results from enumeration of orders 2--200, excluding 64, 96, 128, 144, 160, 162 and 192.
    \item \texttt{results/} --- stores results from enumeration.
    \item \texttt{logs/} --- stores \texttt{failed-groups.log} and \texttt{parent\_run.log}, and temporarily stores the manager and worker log files.
    \item \texttt{reports/} --- contains summary reports from all results files mentioned above.
\end{itemize}

\section*{2. Setup and configuration}

Our code assumes that you use the manager/worker environment provided by Magma; see the following online documentation:

\begin{center}
\url{https://magma.maths.usyd.edu.au/magma/handbook/text/70}
\end{center}

To use this, you need to edit \texttt{manager.m} and \texttt{worker.m}:

\begin{itemize}[noitemsep]
    \item \texttt{NumberOfWorkers} is set in \texttt{manager.m}.
    \item The name of the machine and the port used is set in both \texttt{manager.m} and \texttt{worker.m}. Data stored in both must agree.
\end{itemize}

\section*{3. Running the main script (\texttt{parent})}

The main bash shell script is \texttt{parent}. It runs the main function in \texttt{manager.m} on a range of groups.

Edit \texttt{parent} as follows:

\begin{itemize}[noitemsep]
    \item Set \texttt{ORDERS} to the sequence of orders you wish to process.
    \item Set \texttt{M\_VALUES} to the sequence of \texttt{NumberOfGroups(i)} for those orders \texttt{i} listed in \texttt{ORDERS}.
\end{itemize}

You can generate both sequences using \texttt{generate-arrays.m} and then paste them into \texttt{parent}. Run as follows:

\begin{lstlisting}
% ./parent
\end{lstlisting}

By default, it generates \texttt{Report-*.txt} for each \texttt{*} in the sequence \texttt{ORDERS}, stored in the \texttt{reports} directory. A complete record of the run is recorded in \texttt{logs/parent\_run.log}.

\section*{4. Processing an individual order (\texttt{script})}

To process an individual order, use \texttt{script}. Edit the following parameters:

\begin{lstlisting}
m=1
M=14   # Run from m to M
order=16  # Order parameter to pass to Magma
\end{lstlisting}

This processes groups \texttt{[m..M]} of the given order. You can also set \texttt{TIMEOUT} to a shorter or longer time. If a job runs out of time, it is reported in the file specified by \texttt{FAILURE\_LOG} in \texttt{script}. Run as follows:

\begin{lstlisting}
% ./script
\end{lstlisting}

\section*{5. Generating a summary for an individual order}

To generate a summary for an individual order $n$ after running \texttt{./script}:

\begin{lstlisting}
python3 extract_magma_logs.py
\end{lstlisting}

This collects data from \texttt{logs/manager*.log} and generates \texttt{reports/Report-n.txt}, in the same format as the files currently in \texttt{reports/}.

\section*{6. Processing results}

The data for the groups of order $n$ is recorded in \texttt{results/Result-n}. This can be summarised by editing \texttt{process-results.m} to list the relevant values of $n$, then running:

\begin{lstlisting}
% magma process-results.m 
\end{lstlisting}

This reports the total for each order $n$, where each report corresponds to a row in the enumeration table in the preprint.

Note that \texttt{Report-n.txt} and \texttt{Results-n} are generated in the \texttt{reports} and \texttt{results} directories respectively.

\section*{7. Running interactively}

You can also run a job interactively to enumerate data for an individual group in Magma. Uncomment and edit these parameters in \texttt{manager.m} (retain the quotes):

\begin{lstlisting}
// order := "81";
// nmr := "2";
\end{lstlisting}

Then run:

\begin{lstlisting}
% magma manager.m
\end{lstlisting}

This runs the enumeration algorithm for \texttt{SmallGroup(81, 2)}, reports the outcome, and then exits Magma.

\section*{8. Utility shell scripts}

The utility shell scripts require execute permission and perform the following tasks:

\begin{itemize}[noitemsep]
    \item \texttt{delete-jobs} --- deletes all Magma jobs which were started in the current directory.
    \item \texttt{replace} \textit{old-string} \textit{new-string} --- replaces a string (e.g.\ machine name) in \texttt{manager*.m} and \texttt{worker*.m}.
\end{itemize}

\section*{9. Standard serial run}

For a standard serial run (ignoring parallel option), we provide {\tt 
serial-run.m}.

Edit the values of {\tt order} and {\tt nmr} in that file 
to process {\tt SmallGroup (order, nmr)}.

\end{document}
